\chapter{Babesiosis}
\index{Babesiosis}

\section*{Definition and Overview}
Babesiosis is a rare infectious disease caused by Babesia parasites, microscopic organisms that invade and destroy red blood cells. It is spread to people mainly by the bite of an infected tick, and it is closely related to malaria, which is caused by a similar type of parasite.

The parasites that cause babesiosis are found in different parts of the world, including the United States, parts of Europe and Asia. The tick that carries Babesia is the same species that transmits Lyme disease, so the two infections can occur together. Many people infected with Babesia have no symptoms at all, but others develop a flu-like illness that can become severe, particularly in older people and those with weakened immune systems.

\section*{Epidemiology}
Babesiosis occurs mainly in the north-eastern and mid-western United States, with smaller numbers of cases reported in Europe and elsewhere. Cases follow the distribution of the deer (black-legged) tick and peak in the warmer months from spring to autumn, when ticks are most active. People who spend time outdoors in wooded or grassy areas are at highest risk.

Most infections in healthy people are mild or symptom-free, so the true number of cases is higher than reported. Severe disease is most common in older adults, people who have had their spleen removed, and people with weakened immune systems. Babesia can also be transmitted through blood transfusion and, rarely, from a mother to her baby during pregnancy or birth.

\section*{Aetiology and Pathophysiology}
Babesiosis is caused by protozoan parasites of the Babesia family. The main species affecting people are \textit{Babesia microti} in North America and \textit{Babesia divergens} in Europe.

The infection is usually acquired when an infected tick bites a person and feeds on their blood. The parasite enters the bloodstream and invades red blood cells, where it multiplies. When the infected red cells burst, more parasites are released and further cells are infected. This cycle destroys red blood cells, causing anaemia, and triggers a strong immune and inflammatory response. Severe infection can damage the kidneys, liver and lungs, and can impair the clotting of blood.

\section*{Clinical Features}
Symptoms typically begin one to four weeks after a tick bite, though many infections cause no symptoms at all.

\begin{itemize}
    \item Fever, chills and drenching sweats
    \item Headache and fatigue
    \item Muscle and joint aches
    \item Nausea, vomiting and poor appetite
    \item Cough and shortness of breath, in some cases
    \item Pallor and jaundice (yellowing of the skin and eyes) from destruction of red blood cells
    \item Dark-coloured urine in severe disease
    \item Severe disease may progress to confusion, bleeding and organ failure, mainly in high-risk people
\end{itemize}

\section*{Differential Diagnosis}
The early symptoms of babesiosis are easily confused with other illnesses, especially in areas where the infection is uncommon.

\begin{itemize}
    \item Lyme disease, which is carried by the same tick and can occur at the same time
    \item Anaplasmosis and ehrlichiosis, other tick-borne infections with similar symptoms
    \item Malaria, which is clinically very similar and is diagnosed by the same kind of blood test
    \item Influenza and COVID-19
    \item Glandular fever (infectious mononucleosis)
    \item Sepsis from other causes
\end{itemize}

\section*{When to Seek Medical Care}
See a doctor promptly if you develop fever, chills or flu-like symptoms after a tick bite or after spending time in a tick-infested area, especially if you are older, have no spleen, or have a weakened immune system. Mention the possible tick exposure, as it strongly influences testing.

\textbf{Contact your local emergency services for:}
\begin{itemize}
    \item Difficulty breathing or chest pain
    \item Confusion, drowsiness or fainting
    \item Unusual bruising or bleeding
    \item Very dark urine, or greatly reduced urine output
    \item Persistent high fever that is not improving
\end{itemize}

\section*{Diagnosis and Investigations}
Babesiosis is diagnosed by finding the parasite in the blood or by blood tests that detect its genetic material or antibodies.

\begin{itemize}
    \item \textbf{History and examination:} recent tick exposure, travel and symptoms
    \item \textbf{Blood film:} a sample of blood examined under a microscope to see parasites inside red cells
    \item \textbf{PCR testing:} a sensitive test that detects Babesia DNA in the blood
    \item \textbf{Serology:} blood tests for antibodies, used mainly to confirm past infection
    \item \textbf{Full blood count:} to check for anaemia and a low platelet count
    \item \textbf{Kidney and liver tests:} to detect organ involvement in severe cases
    \item \textbf{Testing for co-infections:} because the same tick can carry Lyme disease and anaplasmosis
\end{itemize}

\section*{Management}
Healthy people with mild infection may not need any treatment, as the illness often resolves on its own. Moderate to severe cases need medicines, and the most severe cases need hospital care.

\begin{itemize}
    \item \textbf{Antiparasitic medicines:} a combination such as atovaquone with azithromycin is commonly used; a different combination is used for severe disease
    \item \textbf{Supportive care:} fluids, rest and medicines for fever and pain
    \item \textbf{Blood transfusion:} for significant anaemia
    \item \textbf{Exchange transfusion:} a procedure that removes infected red cells, used in life-threatening cases
    \item \textbf{Hospital and intensive care:} for people with organ failure or breathing difficulty
    \item \textbf{Monitoring:} repeat blood tests to confirm the parasites are clearing, especially in high-risk people
\end{itemize}

\section*{Complications}
\begin{itemize}
    \item Severe haemolytic anaemia requiring transfusion
    \item Kidney failure
    \item Acute respiratory distress syndrome (severe breathing failure)
    \item Liver dysfunction and jaundice
    \item Blood-clotting failure and abnormal bleeding
    \item Relapse, and persistent low-level infection that can last for months
    \item Death, mainly in the elderly, people without a spleen and people with weakened immunity
\end{itemize}

\section*{Prognosis and Outcomes}
Most healthy people recover fully, and those with mild infection may recover without treatment. Treatment is effective for moderate disease, with most people improving over one to two weeks. The outlook is much less favourable for older adults, people who have had their spleen removed, and people with weakened immune systems, who are at higher risk of severe complications and death. Even after recovery, low levels of parasite may remain in the blood for months, which is why blood donors are screened in affected regions.

\section*{Prevention and Self-Care}
Prevention centres on avoiding tick bites, since there is no vaccine.

\begin{itemize}
    \item Use insect repellent containing DEET or picaridin on exposed skin
    \item Wear long sleeves, long trousers and light-coloured clothing in tick areas
    \item Treat clothing and outdoor gear with permethrin where available
    \item Check yourself, children and pets for ticks after being outdoors, and check again later
    \item Remove attached ticks promptly with fine tweezers, gripping close to the skin and pulling steadily
    \item Shower soon after outdoor activities in tick areas
    \item See a doctor early if you develop fever after a tick bite or time in tick country
    \item Blood donors in affected regions are screened; do not donate blood if you are unwell after a possible tick bite
\end{itemize}

\section*{Sources and Further Reading}
The following reputable organisations provide reliable, current information on babesiosis.

\begin{itemize}
    \item Centers for Disease Control and Prevention. \textit{Information on babesiosis, symptoms and treatment.} https://www.cdc.gov/babesiosis/
    \item National Health Service. \textit{Information on babesiosis and other tick-borne infections.} https://www.nhs.uk/conditions/babesiosis/
    \item MedlinePlus. \textit{Health information on babesiosis.} https://medlineplus.gov/babesiosis.html
    \item World Health Organization. \textit{Information on vector-borne diseases and tick-borne infections.} https://www.who.int/health-topics/vector-borne-diseases
    \item MSD Manual. \textit{Consumer information on babesiosis.} https://www.msdmanuals.com/home/infections
    \item Mayo Clinic. \textit{Information on babesiosis and tick-borne illness.} https://www.mayoclinic.org/diseases-conditions/babesiosis
\end{itemize}
